\lecture{9}{Mon 19 Sep 2022 11:32}{Compactness and Heine-Borel theorem}

Suppose we have $E \subset \mathbb{R}^p$ and a family $\mathscr{G} = \{G_\alpha\} _{\alpha\in A}$ where $G_\alpha \subset \mathbb{R}^p$ is open for every $\alpha \in A$.
We say that $\mathcal{G}$ is an \textit{open cover} of $E$ if \[
	E \subset \bigcup_{\alpha \in A} G_\alpha
.\] 

We say that $\mathcal{G}' \subset \mathcal{G}$ is a subcover if we still have \[
	E \subset \bigcup_{G \in \mathcal{G}'} G
.\] 

\begin{definition}[Compact]
	\label{thm:heine-borel}
	$E \subset \mathbb{R}^p $ is called \textit{compact} if every open cover $\mathcal{G}$ of $E$ has a finite subcover.
\end{definition}

\begin{example}
	\begin{enumerate}
		\item If $E$ is finite, then $E$ is compact.
			\begin{proof}
				Let $E = \{x_1,x_2,\ldots,x_N\} $. Let $\mathcal{G}$ be an open cover. 
				For $x_1, \exists G_1 \in \mathcal{G}$ such that $x_1 \in G_1$.
				For $x_2, \exists G_2 \in \mathcal{G}$ such that $x_2 \in G_2$. Continue in this fashion, we have a finite family of open sets \[
				E \subset \bigcup_{i = 1} ^N G_i
				.\] 
				Hence $\{G_1, \ldots, G_N\} $ is a subcover of $\mathcal{G}$.
			\end{proof}
		\item $\mathbb{R}^p$ is not compact. 
		\begin{proof}
		Let $G_n = B_n(0)$. Then 
		\[
			\mathbb{R}^p \subset \bigcup_{n= 1} ^\infty G_n
		.\] 
		But \begin{align*}
			\mathbb{R}^p &\not \subset G_{n_1} \cup \ldots \cup G_{n_N}\\
			&= B_{n_1}(0) \cup  \ldots \cup B_{n_N}(0)\\
			&= B_{\max n_i}(0)
		.\end{align*} 
			
		\end{proof}
		
	\end{enumerate}
\end{example}

\begin{theorem}[Heine-Borel Theorem]
	$E\subset \mathbb{R}^p$ is compact if and only if $E$ is closed and bounded.
\end{theorem}
\begin{remark}
	True only in $\mathbb{R}^p$, finite-dimensional Euclidean space. Might not be true in spaces other than $\mathbb{R}^p$.
\end{remark}

\begin{lemma}
	If $E$ compact, then $E$ bounded.
\end{lemma}
\begin{proof}
	Take $G_{n} = B_n(0) , n = 1,2,\ldots$
	Then $\{G_n\} $ is an open cover of $E$: \[
		E \subset \bigcup_{n=1} ^\infty G_n = \mathbb{R}^p
	.\] 
	By compactness, there are $G_{n_1}, \ldots, G_{n_N}$ such that \begin{align*}
		E &\subset G_{n_1}\cup \ldots \cup G_{n_N}\\
		&= B_{\max \{n_1, \ldots, n_N\} }(0)
	.\end{align*} 
	Hence $E$ bounded.
\end{proof}
\begin{lemma}
	If $E$ compact, then $E $ closed.
\end{lemma}
\begin{proof}
	Suppose $x$ is a cluster point of $E$ and $x \not\in E$. Let \[
		G_n = \overline{B_{\frac{1}{n}}(x)} ^c = \{y : \|y - x\| >  \frac{1}{n}\} 
	.\] 
	Then \[
		\bigcup_{n = 1} ^\infty G_n = \mathbb{R}^p \setminus \{x\} 
		\supset E
	.\] 
	If there exists a finite subcover $G_{n_1}, \ldots, G_{n_N}$ such that \[
		E \subset G_{n_1} \cup \ldots \cup G_{n_N} = G_{\max \{n_1, \ldots, n_N\} } =: G_R
	.\] 
	Then $B_R(x) \cap E = \emptyset $. This is a contradiction since $x$ is a cluster (or boundary) point. Hence for every cluster point $x$ of $E$ we have $x \in E$.
	Hence $E$ closed.
\end{proof}
\begin{remark}
	The two lemmas above are true in every topological space. It's the reverse statement (closed and bounded implies compactness) that isn't universal.
\end{remark}

\begin{lemma}
	Closed cellls are compact.
\end{lemma}
\begin{proof}
	(Bisection argument) Let $I$ be a closed cell and $\mathcal{G}$ an open cover of $I$.
	Suppose $\mathcal{G}$ does not have finite subcover.
	\begin{enumerate}
		\item Bisect the sides of $I$ to form $2^p$ closed cells. If all $2^p$ subcells have finite subcover of $\mathcal{G}$,  $\mathcal{G}$ would have a finite subcover of $I$.
			Hence there exists a cell, call it $I_1$, such that $I_1$ cannot be covered by finitely many elements of $\mathcal{G}$. We have $l(I_1) = \frac{1}{2}l(I)$.
		\item Bisect $I_1$ to find $I_2$, in the same fashion. By induction, we construct \[
			I_1 \supset I_2 \supset \ldots \supset I_n \supset \ldots
		.\] 
		Such that $l(I_n) = \frac{1}{2^n}l(I)  $, and $I_n$ cannot be covered by finitely many elements of $\mathcal{G}$.
		\item
		Now take $\xi \in \bigcap_{n=1} ^\infty I_n$ (nonempty by Nested Cells Lemma). $\xi \in I$, so there $\exists G \in \mathcal{G}$ suc that $\xi \in G$. Since $G$ is open, there $\exists \delta>0$ such that $B_\delta(\xi)\subset G$.

		As we have shown in the proof of Bolzano-Weierstrass, there exists an $n$ such that $I_n \subset B_\delta(\xi) \subset G$.
		This means $I_n$ is covered by $1$ element of $\mathcal{G}$, contrary to the assumption that $I_n$ cannot be covered by finite subcover of $\mathcal{G}$.
		Hence $I$ is compact.
	\end{enumerate}
\end{proof}

\begin{lemma}
	If $F$ is compact, then $E\subset F$ is compact iff it is closed.
\end{lemma}
\begin{proof}
	Let $\mathcal{G}$ be open cover of $E$. Add $E^c$ to $\mathcal{G}$, and  \[
		\tilde {\mathcal{G}} = \mathcal{G} \cup \{E^c\} 
	.\] 
	Now \[
		\bigcup_{G \in \tilde{\mathcal{G}}} G = \left( \bigcup_{G \in \mathcal{G}} G \right) \cup  E^c = \mathbb{R}^p \supset F
	.\] 
	Hence $\tilde{\mathcal{G}}$ is an open cover of $F$. Hence there exists a finite subcover of $F$.
	
	If $\{G_1, \ldots, G_N\} $ contains $E^c$, throw it out and call the remaining finite set $\mathcal{G}'$. Then $\mathcal{G}'$ covers $E$, and $\mathcal{G}' \subset \mathcal{G}$. Therefore $E$ is compact.
\end{proof}

\begin{proof}[Proof of Theorem \ref{thm:heine-borel}]
	By the four lemmas.
\end{proof}

